\begin{Exercise}[label={ex:espr_logica}]
Si consideri la seguente espressione booleana:
$$((!A \ || \ B) \ \&\& \ (A \  \&\& !B))\ || \ C$$
e se ne compili la corrispondente tabella di verit\`a (ossia una tabella in cui $0$ rappresenta il valore logico FALSO, $1$ il valore VERO).

Si consideri ora la condizione, scritta in linguaggio C, in cui x e y siano due variabili intere:
$$((!(x < 0) || (y>0)) \ \&\& \ ((x < 0) \ \&\& \ !(y > 0))) ) \ || \ (x > 3)$$
ottenuta dall'espressione booleana sopra indicata sostituendo la variabile $A$ con $x < 0$, la variabile $B$ con $y > 0$, la variabile $C$ con $x > 3$
Si risponda alle seguenti domande giustificando la risposta:
\begin{enumerate}
	\item L'espressione \`e vera o falsa quando $x = 5$ e $y = 7$?
	\item Se $x < 0$, per quali valori di y l'espressione \`e vera?
\end{enumerate}
\end{Exercise}

\begin{Answer}[ref={ex:espr_logica}]
\begin{tabular}{|c|c|c|c|c|c|c|c|c|}
\hline
A & B & C & !A & !B & (!A \ || \ B) & (A \  \&\& !B)) & ((!A \ || \ B) \ \&\& \ (A \  \&\& !B)) & TOT\\
\hline
0 & 0 & 0 & 1 & 1 & 1 & 0 & 0 & 0\\
\hline
0 & 0 & 1 & 1 & 1 & 1 & 0 & 0 & 1\\
\hline
0 & 1 & 0 & 1 & 0 & 1 & 0 & 0 & 0\\
\hline
0 & 1 & 1 & 1 & 0 & 1 & 0 & 0 & 1\\
\hline
1 & 0 & 0 & 0 & 1 & 0 & 1 & 0 & 0\\
\hline
1 & 0 & 1 & 0 & 1 & 0 & 1 & 0 & 1\\
\hline
1 & 1 & 0 & 0 & 0 & 1 & 0 & 0 & 0\\
\hline
1 & 1 & 1 & 0 & 0 & 1 & 0 & 0 & 1\\
\hline
\end{tabular}

Quando $x = 5$, la variabile $C = 1$, quindi l'espressione \`e vera. Se $x < 0$, la variabile $C = 0$, quindi l'espressione \`e falsa per qualsiasi valore di $y$.
\end{Answer}